% SPDX-FileCopyrightText: 2026 Lari Natri <lari.natri@iki.fi>
% SPDX-License-Identifier: GPL-3.0-or-later
%
% Annotation scanner, lyric anchoring, and raw chord memory.
% This file is part of the 'ulsbs' package.
%
% Private module loaded by ulsbs.sty after song/layout/verse and common notation
% helpers. Music loads next and supplies annotation ink; multikey and diagrams
% subsequently adapt the rendering seams defined here. These forward references
% are used at song-render time, not while this file is being loaded.
%
% The scanner collects unexpanded lyric tokens, typesets each fragment into one
% retained box, and consumes that box for metrics and output. Never measure by
% evaluating the lyric source again: it may contain writes or other side effects.
% Music replaces the single-annotation box; this module concatenates it into a
% run and owns whitespace minima, lyric boundaries, alignment, and break costs.
%
% Register declarations and raw memories are global, as are replay consumption
% and the selected register, so they survive nested lyric/typesetting groups.
% Song begin clears only the default memory/current replay, not named registers.
% Stanza start selects default recording or replay; local tracking/scanner state
% follows the surrounding stanza scope. Raw replay is rendered with current
% visibility and transposition, rather than reusing previously rendered ink.

\ProvidesFile{ulsbs-chords.code.tex}[Chord and lyric scanner]

\makeatletter
\ExplSyntaxOn
\box_new:N \l__ulsbs_chord_box
\box_new:N \l__ulsbs_chord_run_box
\box_new:N \l__ulsbs_chord_lyric_box
\tl_new:N \l__ulsbs_chord_lyric_tl
\tl_new:N \l__ulsbs_chord_before_hyphen_tl
\tl_new:N \g__ulsbs_chord_register_tl
\tl_new:N \l__ulsbs_chord_raw_tl
\tl_new:N \l__ulsbs_chord_saved_tl
\tl_new:N \l__ulsbs_chord_saved_melody_tl
\tl_new:N \l__ulsbs_chord_saved_tail_tl
\tl_new:N \l__ulsbs_chord_new_melody_tl
\tl_new:N \l__ulsbs_chord_new_tail_tl
\seq_new:N \g__ulsbs_chord_replay_seq
\prop_new:N \g__ulsbs_chord_memory_prop
\prop_new:N \g__ulsbs_chord_declared_prop
\prop_new:N \g__ulsbs_chord_lyric_chars_prop
\prop_gput:Nnn \g__ulsbs_chord_declared_prop {}{}
\prop_gput:Nnn \g__ulsbs_chord_memory_prop {}{}
\bool_new:N \l__ulsbs_chord_tracking_bool
\bool_new:N \l__ulsbs_chord_word_end_bool
\bool_new:N \l__ulsbs_chord_broken_bool
\bool_new:N \l__ulsbs_chord_adjacent_bool
\bool_new:N \l__ulsbs_chord_excluded_bool
\bool_new:N \l__ulsbs_chord_full_clearance_bool
\bool_new:N \l__ulsbs_chord_empty_anchor_bool
\int_new:N \l__ulsbs_chord_hyphens_int
\int_new:N \g__ulsbs_chord_spacefactor_int
\dim_new:N \l__ulsbs_chord_width_dim
\newattribute\__ulsbs_chord_annotation_attribute
\int_const:Nn \c__ulsbs_chord_annotation_attribute_int {\allocationnumber}
\newattribute\__ulsbs_chord_row_attribute
\int_const:Nn \c__ulsbs_chord_row_attribute_int {\allocationnumber}
\msg_new:nnn {ulsbs-chords}{state}{Invalid~chord~state:~#1.}

% #1 is diagnostic detail; malformed lifecycle/register use is fatal.
\cs_new_protected:Npn \__ulsbs_chord_fail:n #1
  {\msg_fatal:nnn {ulsbs-chords}{state}{#1}}

% #1 names a global raw-token register; the empty default register exists above.
% Declaring an existing name is an error, rather than silently losing its memory.
\DeclareDocumentCommand \newchords {m}
  {\prop_if_in:NnT \g__ulsbs_chord_declared_prop {#1}
     {\__ulsbs_chord_fail:n {duplicate~register~'#1'}}
   \prop_gput:Nnn \g__ulsbs_chord_declared_prop {#1}{}
   \prop_gput:Nnn \g__ulsbs_chord_memory_prop {#1}{}}

% Validate #1 before globally selecting it; selection does not rewind replay.
\cs_new_protected:Npn \__ulsbs_chord_register:n #1
  {\prop_if_in:NnF \g__ulsbs_chord_declared_prop {#1}
     {\__ulsbs_chord_fail:n {unknown~register~'#1'}}
   \tl_gset:Nn \g__ulsbs_chord_register_tl {#1}}

% [name] defaults to the empty register: clear its contents and record new raw
% annotations. Tracking is local even though stored contents are global.
\DeclareDocumentCommand \memorize {O{}}
  {\__ulsbs_chord_register:n {#1}
   \prop_gput:Nnn \g__ulsbs_chord_memory_prop {#1}{}
   \bool_set_true:N \l__ulsbs_chord_tracking_bool}

% Optional [name] selects a register; omission keeps the current one. Split the
% stored separator-prefixed stream and discard its leading empty item to rewind.
% The separator is a token delimiter, never an instruction to execute.
\DeclareDocumentCommand \replay {o}
  {\IfNoValueF{#1}{\__ulsbs_chord_register:n {#1}}
   \bool_set_false:N \l__ulsbs_chord_tracking_bool
   \prop_get:NVN \g__ulsbs_chord_memory_prop \g__ulsbs_chord_register_tl \l_tmpa_tl
   \seq_gset_split:NnV \g__ulsbs_chord_replay_seq {\__ulsbs_chord_separator:}\l_tmpa_tl
   \seq_gpop_left:NN \g__ulsbs_chord_replay_seq \l_tmpa_tl}

% Song-begin hook: reset default memory/replay and alternate-set color sequence.
\cs_new_protected:Npn \__ulsbs_chord_reset:
  {\prop_gput:Nnn \g__ulsbs_chord_memory_prop {}{}
   \tl_gclear:N \g__ulsbs_chord_register_tl
   \seq_gclear:N \g__ulsbs_chord_replay_seq
   \ulsbs@setaltchordscolor{1}}

% Stanza-start hook: an empty default memory records, otherwise replay it from
% the beginning. Named-register use within a stanza is an explicit author choice.
\cs_new_protected:Npn \__ulsbs_chord_stanza:
  {\prop_get:NnN \g__ulsbs_chord_memory_prop {}\l_tmpa_tl
   \tl_if_empty:NTF \l_tmpa_tl {\memorize}{\replay[]}}

% #1 is an output tl variable. Consume globally so nested annotation boxes do
% not restore a consumed chord; exhaustion is an authoring error.
\cs_new_protected:Npn \__ulsbs_chord_pop:N #1
  {\seq_if_empty:NT \g__ulsbs_chord_replay_seq
     {\__ulsbs_chord_fail:n {replayed~chord~has~no~matching~memorized~chord}}
   \seq_gpop_left:NN \g__ulsbs_chord_replay_seq #1}

% Replay always consumes raw memory, even when all annotation ink is hidden.
\cs_new_protected:Npn \__ulsbs_chord_replay:
  {\__ulsbs_chord_pop:N \l__ulsbs_chord_raw_tl
   \exp_args:NV \__ulsbs_chord_begin:n \l__ulsbs_chord_raw_tl}

% Read annotation source only after installing notation catcodes and chordlocals.
% The delimited reader stops at the first ungrouped ]; braces protect inner ].
% End the read group before rendering so scanner state lives in the stanza.
\cs_set_protected:Npn \[
  {\group_begin:
   \char_set_catcode_active:n {35}\char_set_catcode_active:n {38}
   \char_set_catcode_other:n {58}\ulChordLocals
   \__ulsbs_chord_read:w}

% #1 is captured raw source. Memory and lyric scanning are independent of ink;
% the music renderer must suppress hidden payloads without evaluating them.
\cs_new_protected:Npn \__ulsbs_chord_read:w #1]
  {\group_end:\__ulsbs_chord_begin:n {#1}}

% #1 is one raw annotation. Start an independent lyric anchor, clear run state,
% and require an active song/stanza before any ink or lyric scanning occurs.
\cs_new_protected:Npn \__ulsbs_chord_begin:n #1
  {\__ulsbs_song_require:
   \bool_if:NF \g__ulsbs_stanza_active_bool
     {\__ulsbs_chord_fail:n {annotation~outside~a~stanza}}
   \box_clear:N \l__ulsbs_chord_run_box
   \bool_set_false:N \l__ulsbs_chord_full_clearance_bool
   \tl_clear:N \l__ulsbs_chord_lyric_tl
   \int_zero:N \l__ulsbs_chord_hyphens_int
   \bool_set_false:N \l__ulsbs_chord_adjacent_bool
   \__ulsbs_chord_add:n {#1}}

% #1 joins this anchor's annotation run. Leading ^ excludes recording, or merges
% one replayed annotation with the new tail when replaying. Store source tokens
% before rendering; the music layer must not own memory or rescan lyric source.
\cs_new_protected:Npn \__ulsbs_chord_add:n #1
  {\tl_set:Nn \l__ulsbs_chord_raw_tl {#1}
    \bool_set_false:N \l__ulsbs_chord_excluded_bool
    \tl_if_head_eq_charcode:nNTF {#1}^
     {\bool_set_true:N \l__ulsbs_chord_excluded_bool
      \tl_set:Nx \l__ulsbs_chord_raw_tl {\tl_tail:n {#1}}
      \bool_if:NF \l__ulsbs_chord_tracking_bool {\__ulsbs_chord_embedded_replay:}}
     {\bool_if:NT \l__ulsbs_chord_tracking_bool
       {\prop_get:NVN \g__ulsbs_chord_memory_prop \g__ulsbs_chord_register_tl \l_tmpa_tl
        \tl_put_right:Nn \l_tmpa_tl {\__ulsbs_chord_separator: #1}
        \prop_gput:NVV \g__ulsbs_chord_memory_prop \g__ulsbs_chord_register_tl \l_tmpa_tl}}
   \exp_args:NV \__ulsbs_music_annotation:n \l__ulsbs_chord_raw_tl
   \__ulsbs_chord_check_space:}

% Split raw #1 into output tls #2 (optional leading <melody>) and #3 (tail).
% No expansion or search past the head: only a leading block has melody syntax.
\cs_new_protected:Npn \__ulsbs_chord_split:nNN #1#2#3
  {\tl_clear:N #2\tl_set:Nn #3 {#1}
   \tl_if_head_eq_charcode:nNT {#1}<
     {\__ulsbs_chord_split:w #1\q_stop #2#3}}

% Delimited worker: #1 melody content, #2 tail, #3/#4 destination tl variables.
% Preserve the angle delimiters in the saved melody for later annotation parsing.
\cs_new_protected:Npn \__ulsbs_chord_split:w <#1>#2\q_stop #3#4
  {\tl_set:Nn #3 {<#1>}\tl_set:Nn #4 {#2}}

% Consume exactly one saved annotation without starting another lyric scan.
% Saved melody wins; otherwise retain the new melody. Append the new chord tail
% to the saved tail, protecting all source tokens from the assembly expansion.
\cs_new_protected:Npn \__ulsbs_chord_embedded_replay:
  {\__ulsbs_chord_pop:N \l__ulsbs_chord_saved_tl
   \exp_args:NV \__ulsbs_chord_split:nNN \l__ulsbs_chord_raw_tl
     \l__ulsbs_chord_new_melody_tl \l__ulsbs_chord_new_tail_tl
   \exp_args:NV \__ulsbs_chord_split:nNN \l__ulsbs_chord_saved_tl
     \l__ulsbs_chord_saved_melody_tl \l__ulsbs_chord_saved_tail_tl
   \tl_if_empty:NT \l__ulsbs_chord_saved_melody_tl
     {\tl_set_eq:NN \l__ulsbs_chord_saved_melody_tl \l__ulsbs_chord_new_melody_tl}
   \tl_set:Nx \l__ulsbs_chord_raw_tl
     {\exp_not:V \l__ulsbs_chord_saved_melody_tl
      \exp_not:V \l__ulsbs_chord_saved_tail_tl
      \exp_not:V \l__ulsbs_chord_new_tail_tl}}

% Literal following whitespace grants this visible annotation a minimum advance.
% Do this per annotation, before concatenation; hidden/zero-advance ink opts out.
\cs_new_protected:Npn \__ulsbs_chord_check_space:
  {\peek_catcode:NTF \c_space_token
     {\ifulsbs@annotationminimumdisabled\else\ifulsbs@annotationvisible
       \hbox_set:Nn \l_tmpa_box {{\ulLyricFontStyle\hspace{\ulChordWhitespaceMinimumWidth}}}
       \dim_compare:nNnT {\box_wd:N \l__ulsbs_chord_box}<{\box_wd:N \l_tmpa_box}
         {\box_set_wd:Nn \l__ulsbs_chord_box {\box_wd:N \l_tmpa_box}}
     \fi\fi\__ulsbs_chord_append:}
     {\__ulsbs_chord_append:}}

% Move the new box into the run, preserving already-typeset nodes, then look
% ahead without expanding or consuming the next lyric/control token.
\cs_new_protected:Npn \__ulsbs_chord_append:
   {\ifulsbs@annotationfullclearance
      \bool_set_true:N \l__ulsbs_chord_full_clearance_bool\fi
    % Tag retained ink too: no-lyrics output can unpack the run without a stack.
    \directlua{
      local~box=tex.box[\number\l__ulsbs_chord_box]
      local~attribute=\int_use:N\c__ulsbs_chord_annotation_attribute_int
      node.set_attribute(box,attribute,1)}
   \hbox_set:Nn \l__ulsbs_chord_run_box
     {\hbox_unpack_drop:N \l__ulsbs_chord_run_box\box_use_drop:N \l__ulsbs_chord_box}
   \bool_set_true:N \l__ulsbs_chord_adjacent_bool
   \__ulsbs_chord_scan:}

% Refresh l_peek_token before dispatch; no expansion means source effects remain
% deferred to the single lyric box or to commands after this fragment ends.
\cs_new_protected:Npn \__ulsbs_chord_scan:
  {\peek_after:Nw \__ulsbs_chord_dispatch:}

% Adjacent annotations/replay share a run only before lyrics begin. Quote tokens
% in that position are emitted directly and scanning resumes. All other input
% goes to the lyric classifier; a later annotation ends the current fragment.
\cs_new_protected:Npn \__ulsbs_chord_dispatch:
  {\token_if_eq_meaning:NNTF \l_peek_token \[
     {\bool_if:NTF \l__ulsbs_chord_adjacent_bool
       {\__ulsbs_chord_adjacent:w}{\__ulsbs_chord_emit:}}
     {\bool_if:NTF \l__ulsbs_chord_adjacent_bool
       {\token_if_eq_meaning:NNTF \l_peek_token \__ulsbs_chord_replay:
         {\__ulsbs_chord_adjacent_replay:N}
         {\peek_charcode:NTF ` {\__ulsbs_chord_quote:N}
           {\peek_charcode:NTF ' {\__ulsbs_chord_quote:N}
             {\peek_charcode:NTF " {\__ulsbs_chord_quote:N}
               {\__ulsbs_chord_dispatch_lyric:}}}}}
       {\__ulsbs_chord_dispatch_lyric:}}}

% #1 is one literal quote already recognized by dispatch; leave it outside the
% lyric box and preserve the adjacent-annotation phase while scanning onward.
\cs_new_protected:Npn \__ulsbs_chord_quote:N #1 {#1\__ulsbs_chord_scan:}

% #1 consumes the replay token already recognized by lookahead, not its meaning.
\cs_new_protected:Npn \__ulsbs_chord_adjacent_replay:N #1
  {\__ulsbs_chord_pop:N \l__ulsbs_chord_raw_tl
   \exp_args:NV \__ulsbs_chord_add:n \l__ulsbs_chord_raw_tl}

% Consume adjacent \[ and read its ]-delimited payload under chordlocals. The
% stanza has already installed accidental catcodes for this continuing run.
\cs_new_protected:Npn \__ulsbs_chord_adjacent:w \[
  {\group_begin:\ulChordLocals\__ulsbs_chord_adjacent_read:w}

% #1 is the adjacent raw payload; end capture scope before appending to this run.
\cs_new_protected:Npn \__ulsbs_chord_adjacent_read:w #1]
  {\group_end:\__ulsbs_chord_add:n {#1}}

% Groups are opaque lyric units; spaces stop without being consumed. Distinguish
% control/active tokens from characters without expanding either kind.
\cs_new_protected:Npn \__ulsbs_chord_dispatch_lyric:
  {\bool_set_false:N \l__ulsbs_chord_adjacent_bool
   \peek_catcode:NTF \c_group_begin_token {\__ulsbs_chord_group:n}
     {\peek_catcode:NTF \c_group_end_token {\__ulsbs_chord_group_end:}
       {\peek_catcode:NTF \c_space_token {\__ulsbs_chord_emit:}
         {\ifcat\noexpand\l_peek_token\relax
            \cs_set_eq:NN \__ulsbs_chord_next: \__ulsbs_chord_command:N
          \else
            \cs_set_eq:NN \__ulsbs_chord_next: \__ulsbs_chord_character:
          \fi\__ulsbs_chord_next:}}}}

% #1 is a control token: declared lyric commands stay unexpanded in the fragment;
% any other command runs only after emission, with its following arguments intact.
\cs_new_protected:Npn \__ulsbs_chord_command:N #1
  {\prop_if_in:NnTF \g__ulsbs_chord_lyric_chars_prop {#1}
     {\tl_put_right:Nn \l__ulsbs_chord_lyric_tl {#1}\__ulsbs_chord_scan:}
     {\__ulsbs_chord_emit: #1}}

% #1 is a balanced group's contents. Restore its braces rather than flattening
% local formatting or inspecting argument tokens as standalone lyric characters.
\cs_new_protected:Npn \__ulsbs_chord_group:n #1
  {\tl_put_right:Nn \l__ulsbs_chord_lyric_tl {{#1}}\__ulsbs_chord_scan:}

% A closing brace cannot be taken as an ordinary undelimited macro argument.
% The incomplete let assignment consumes it via afterassignment lookahead; the
% next token decides word continuation, then emission occurs before the explicit
% egroup restores the authored scope. Do not replace this with early group end:
% the final syllable still needs the group's font/state during typesetting.
\cs_new_protected:Npn \__ulsbs_chord_group_end:
   {\tex_afterassignment:D \__ulsbs_chord_group_end_scan:
    \cs_set_eq:NN \__ulsbs_chord_next:}

\cs_new_protected:Npn \__ulsbs_chord_group_end_scan:
   {\peek_after:Nw \__ulsbs_chord_group_end_emit:}

\cs_new_protected:Npn \__ulsbs_chord_group_end_emit:
   {\__ulsbs_chord_emit:\egroup}

% #1 is one ordinary lyric token, appended unexpanded before the next lookahead.
\cs_new_protected:Npn \__ulsbs_chord_take:N #1
  {\tl_put_right:Nn \l__ulsbs_chord_lyric_tl {#1}\__ulsbs_chord_scan:}

% Remember the pre-hyphen fragment once, gather the entire literal hyphen run,
% and emit before the first following character. One hyphen requests alignment;
% multiple hyphens remain literal text and suppress synthetic extension.
\cs_new_protected:Npn \__ulsbs_chord_character:
  {\peek_charcode:NTF -
     {\int_compare:nNnT {\l__ulsbs_chord_hyphens_int}={0}
        {\tl_set_eq:NN \l__ulsbs_chord_before_hyphen_tl \l__ulsbs_chord_lyric_tl}
      \int_incr:N \l__ulsbs_chord_hyphens_int
      \__ulsbs_chord_take_hyphen:N}
     {\int_compare:nNnTF {\l__ulsbs_chord_hyphens_int}>{0}
       {\__ulsbs_chord_emit:}{\__ulsbs_chord_take:N}}}

% #1 consumes the hyphen counted above; keep it literal until emit chooses whether
% to replace a single hyphen or retain a multi-hyphen run.
\cs_new_protected:Npn \__ulsbs_chord_take_hyphen:N #1
  {\tl_put_right:Nn \l__ulsbs_chord_lyric_tl {#1}\__ulsbs_chord_scan:}

% #1 is the control token allowed to stay in lyric source. This declaration does
% not evaluate it or parse its arguments; balanced arguments remain groups above.
\DeclareDocumentCommand \ulDeclareLyricChar {m}
  {\prop_gput:Nnn \g__ulsbs_chord_lyric_chars_prop {#1}{}}
\clist_map_inline:nn {\`,\',\^,\",\~,\=,\.,\u,\v,\H,\t,\c,\d,\b,
  \oe,\OE,\ae,\AE,\aa,\AA,\o,\O,\l,\L,\ss,\i,\j,\/,\-,\discretionary}
  {\ulDeclareLyricChar{#1}}

% Finish the fragment using the nonexpanded lookahead token as boundary context.
% The punctuation test is by catcode (ordinary-character category), not a list
% of punctuation glyphs. A single source hyphen is replaced by alignment ink.
% Typeset lyric tokens exactly once here, retain spacefactor across the stack,
% and let Lua fit alternate chord nodes against this retained lyric width.
% All later dimension/break decisions inspect boxes, never reevaluate source.
\cs_new_protected:Npn \__ulsbs_chord_emit:
  {\bool_set_false:N \l__ulsbs_chord_word_end_bool
   \token_if_eq_catcode:NNT \l_peek_token \c_space_token
     {\bool_set_true:N \l__ulsbs_chord_word_end_bool}
    \token_if_eq_catcode:NNT \l_peek_token \c_group_end_token
      {\bool_set_true:N \l__ulsbs_chord_word_end_bool}
    \token_if_eq_meaning:NNT \l_peek_token \endgroup
      {\bool_set_true:N \l__ulsbs_chord_word_end_bool}
    \token_if_eq_meaning:NNT \l_peek_token \par
      {\bool_set_true:N \l__ulsbs_chord_word_end_bool}
    \token_if_eq_meaning:NNT \l_peek_token \__ulsbs_verse_line:
      {\bool_set_true:N \l__ulsbs_chord_word_end_bool}
   \token_if_eq_catcode:NNT \l_peek_token .
     {\bool_set_true:N \l__ulsbs_chord_word_end_bool}
   \int_compare:nNnT {\l__ulsbs_chord_hyphens_int}={1}
     {\tl_set_eq:NN \l__ulsbs_chord_lyric_tl \l__ulsbs_chord_before_hyphen_tl}
    \hbox_set:Nn \l__ulsbs_chord_lyric_box
      {\tl_use:N \l__ulsbs_chord_lyric_tl
       \int_gset:Nn \g__ulsbs_chord_spacefactor_int {\spacefactor}}
   \bool_set:Nn \l__ulsbs_chord_empty_anchor_bool
     {\dim_compare_p:nNn
       {\box_wd:N \l__ulsbs_chord_run_box+\box_ht:N \l__ulsbs_chord_run_box
         +\box_dp:N \l__ulsbs_chord_run_box+\box_wd:N \l__ulsbs_chord_lyric_box}={0pt}}
   \ifulsbs@pdflayers
     \directlua{ulsbs_multikey.fit_chords(\number\l__ulsbs_chord_run_box,
       \number\wd\l__ulsbs_chord_lyric_box,\ulPdfTranspositionMinimumScale)}
   \fi
   \bool_set_false:N \l__ulsbs_chord_broken_bool
   \dim_compare:nNnT {\box_wd:N \l__ulsbs_chord_lyric_box}>{0pt}
     {\dim_compare:nNnT {\box_wd:N \l__ulsbs_chord_run_box}>{\box_wd:N \l__ulsbs_chord_lyric_box}
       {\bool_if:NF \l__ulsbs_chord_word_end_bool
         {\bool_set_true:N \l__ulsbs_chord_broken_bool}}}
   \int_compare:nNnTF {\l__ulsbs_chord_hyphens_int}={1}
     {\bool_set_true:N \l__ulsbs_chord_broken_bool}
     {\int_compare:nNnT {\l__ulsbs_chord_hyphens_int}>{1}
       {\bool_set_false:N \l__ulsbs_chord_broken_bool}}
   \leavevmode
    \legacy_if:nTF {ulNoLyrics}
      {\dim_compare:nNnTF {\box_wd:N \l__ulsbs_chord_lyric_box}={0pt}
        {% Keep beat-only ink in its lane without changing chord-only rows.
         \bool_if:NTF \l__ulsbs_chord_full_clearance_bool
          {\hbox_unpack_drop:N \l__ulsbs_chord_run_box}
          {\__ulsbs_chord_stack:}}
       {\__ulsbs_chord_stack:}}
     {\__ulsbs_chord_stack:}
   \bool_if:NF \l__ulsbs_chord_word_end_bool
     {\int_compare:nNnTF {\l__ulsbs_chord_hyphens_int}>{0}
       {\penalty\exhyphenpenalty}
       {\bool_if:NTF \l__ulsbs_chord_broken_bool {\penalty\hyphenpenalty}{\nobreak}}
       \spacefactor=\g__ulsbs_chord_spacefactor_int\relax}
    % Collapse two source spaces across hidden ink, but preserve a sole following
    % space (including after an empty group or an active measure bar).
    \bool_if:NT \l__ulsbs_chord_empty_anchor_bool
      {\int_compare:nNnT {\tex_lastnodetype:D}={11}{\ignorespaces}}}

% Consume run/lyric boxes into a two-line stack. A broken syllable reserves the
% wider line plus extension space and uses the current font's hyphen character.
% An empty lyric anchor preserves preceding glue but adds a modest break cost.
\cs_new_protected:Npn \__ulsbs_chord_stack:
  {\dim_compare:nNnTF
     {\box_wd:N \l__ulsbs_chord_run_box+\box_ht:N \l__ulsbs_chord_run_box
       +\box_dp:N \l__ulsbs_chord_run_box}={0pt}
     {\hbox_unpack_drop:N \l__ulsbs_chord_lyric_box
      \int_compare:nNnT {\l__ulsbs_chord_hyphens_int}={1}{-}}
     {\__ulsbs_chord_stack_visible:}}

% Visible retained extents alone require a stack; an empty run must not reserve
% the nominal annotation-to-lyric baseline or a blank first hbox.
\cs_new_protected:Npn \__ulsbs_chord_stack_visible:
   {\dim_compare:nNnT {\box_wd:N \l__ulsbs_chord_lyric_box}={0pt}
     {\skip_set:Nn \l_tmpa_skip {\tex_lastskip:D}\unskip\penalty200\hskip\l_tmpa_skip}
    \vbox~attr~\c__ulsbs_chord_annotation_attribute_int~1~
     {\bool_if:NTF \l__ulsbs_chord_full_clearance_bool
       {\ulsbs@annotationfullclearancetrue}{\ulsbs@annotationfullclearancefalse}
     \ulChordLineParams
     \bool_if:NTF \l__ulsbs_chord_broken_bool
       {\dim_set:Nn \l__ulsbs_chord_width_dim
          {\dim_max:nn {\box_wd:N \l__ulsbs_chord_run_box}{\box_wd:N \l__ulsbs_chord_lyric_box}+.5em}
        \hbox_to_wd:nn {\l__ulsbs_chord_width_dim}{\hbox_unpack_drop:N \l__ulsbs_chord_run_box\hfil}
        \hbox_to_wd:nn {\l__ulsbs_chord_width_dim}
          {\hbox_unpack_drop:N \l__ulsbs_chord_lyric_box\hfil
           \ifnum\hyphenchar\font>-1\char\hyphenchar\font\hfil\fi}}
       {\box_use_drop:N \l__ulsbs_chord_run_box
        \hbox{\hbox_unpack_drop:N \l__ulsbs_chord_lyric_box\hfil}}}}

% #1 is alternate-set number, #2 chord text. Inline stanza references participate
% in PDF layers; standalone alternate sets use their independent source policy.
% Inline repetition references are not standalone sets: altchords-off keeps them.
\RenewDocumentCommand \ac {m m}
  {\legacy_if:nT {ulShowChords}{\group_begin:
   \str_if_eq:nnTF {#1}{1}{\colorlet{ulsbs@accolor}{.}}
     {\bool_if:NTF \l__ulsbs_chord_in_altchords_bool
       {\colorlet{ulsbs@accolor}{.!50!ulAltInlineChordColor}}
       {\colorlet{ulsbs@accolor}{ulAltInlineChordColor}}\color{ulsbs@accolor}}
   \textsuperscript{#1}
    \bool_if:NTF \l__ulsbs_chord_in_altchords_bool
      {\color{ulsbs@accolor}\ulsbs@transposealtchord{#2}}
      {\bool_if:NT \g__ulsbs_stanza_active_bool {\ulsbs@pdfchordtrue}
       \ulChordFont\color{ulsbs@accolor}\transposehere{#2}}
    \group_end:}}

% #1 transposes only on opt-in; suppress song-key spelling for standalone sets.
\cs_set_protected:Npn \ulsbs@transposealtchord #1
  {\ifulTransposeAltChords\group_begin:\ulsbs@keytransposefalse\transposehere{#1}\group_end:
   \else#1\fi}

% Install accidental catcodes before reading the alternate-set argument, then
% render only when both chord and alternate-set visibility permit it.
\cs_set_protected:Npn \altchords
  {\group_begin:\char_set_catcode_active:n {35}\char_set_catcode_active:n {38}
   \__ulsbs_chord_alternates:n}

% #1 is the captured alternate-set source. Local id, meterfont, and e definitions
% adapt common formatting; disable PDF chord/key policy for this standalone set.
\cs_new_protected:Npn \__ulsbs_chord_alternates:n #1
  {\group_end:
   \bool_if:nT {\legacy_if_p:n {ulShowChords} && \legacy_if_p:n {ulShowAltChords}}
   {\group_begin:
     \bool_set_true:N \l__ulsbs_chord_in_altchords_bool\ulsbs@pdfchordfalse\ulsbs@keytransposefalse
     \cs_set_protected:Npn \id {\@ifnextchar[\__ulsbs_chord_alt_id:w
       {\__ulsbs_chord_alt_id:w [1]}}
     \cs_set:Npn \ulMeterFont {\color{ulChordColor}\normalfont\tiny\sffamily\upshape}
      \cs_set:Npn \e {\ifulShowMeasures-\fi}
     \vspaceifnotinverse{.618034\baselineskip}\hfill
     \ulPrintChord{\raisebox{.5ex}{\ulAltChordsStyle\color{ulsbs@currentaltchordcolor}\ulsbs@transposealtchord{#1}}}
    \group_end:}}

% [#1] is an alternate-set identifier and #2 its chord; reuse common ID styling.
\cs_new_protected:Npn \__ulsbs_chord_alt_id:w [#1]#2
  {\ulsbs@altchordsid[#1]{\ulsbs@transposealtchord{#2}}}
\AddToHook{ulsbs/song/begin}{\__ulsbs_chord_reset:}
\AddToHook{ulsbs/stanza/start}{\__ulsbs_chord_stanza:}
\ExplSyntaxOff
\makeatother
